% from ega3-7-fr.tex
\oldpage[III]{204}
\begin{proof}
La première assertion est conséquence immédiate de (7.6.1) et du fait que la
question est locale sur $Y$ et $Y'$, d'après (7.7.5, \emph{c)} et \emph{b)}). Pour démontrer
la seconde assertion, il suffit de voir que pour tout $y\in Y$,
$T_p^{\mathscr O_y}$ est exact à droite (resp. à gauche), en vertu de (7.7.10).
Mais par hypothèse, il existe $y'\in Y'$ tel que $g(y')=y$, et
$\mathscr O_{y'}$ est un $\mathscr O_y$-module fidèlement plat ; la conclusion
résulte alors de l'hypothèse et de (7.6.1).
\end{proof}

\begin{remarks}[7.7.12]
\textnormal{(i)} Sous les hypothèses de (7.7.4), supposons en outre que
$\mathscr P_\bullet$ soit un complexe \emph{fini} ; alors il résulte de (7.7.1)
(puisqu'on peut prendre pour $\mathfrak U$ un recouvrement ouvert affine
\emph{fini} de $X$) que le bicomplexe
$C^\bullet(\mathfrak U,\mathscr P^\bullet\otimes_{\mathscr O_Y}\mathscr M)$ est
aussi \emph{fini}, et de façon précise qu'il existe un ensemble fini $E$ de couples,
\emph{indépendant de $\mathscr M$}, tel que
$C^h(\mathfrak U,\mathscr P^k\otimes_{\mathscr O_Y}\mathscr M)=0$ pour tous les couples $(h,k)\notin E$.
On en conclut qu'il existe $i_1$ tel que, pour $i\geq i_1$, on ait
$\mathscr T_i(\mathscr M)=0$ pour \emph{tout} $\mathscr O_Y$-Module
quasi-cohérent $\mathscr M$. En particulier, $\mathscr T_i$ est trivialement un
foncteur exact en $\mathscr M$ pour ces valeurs de $i$, et par suite (7.4.1),
$Z'_i(L_\bullet)$ est un $A$-module \emph{plat} de type fini (donc
\emph{projectif} de type fini, puisque $A$ est noethérien) pour ces valeurs de
$i$. Considérons alors le complexe $(L'_\bullet)$ de $A$-modules tel que
$L'_i=L_i$ pour $i<i_1$, $L'_{i_1}=Z'_{i_1}(L_\bullet)$ et $L'_i=0$ pour
$i>i_1$ et soit $\mathscr L'_\bullet=\widetilde{L}'_\bullet$. Il est clair que
$  \mathscr H_i(\mathscr L'_\bullet\otimes_{\mathscr O_Y}\mathscr M)
  =
  \mathscr H_i(\mathscr L'_\bullet\otimes_{\mathscr O_Y}\mathscr M)$
pour $i<i_1-1$ et aussi pour $i\geq i_1$ (les deux membres étant alors nuls) ;
enfin, comme
$\operatorname{Im}(Z'_{i_1}\otimes_A M)=\operatorname{Im}(L_{i_1}\otimes_A M)$
par définition, on a aussi
$\mathscr H_i(\mathscr L'_\bullet\otimes_{\mathscr O_Y}\mathscr M)
=\mathscr H_i(\mathscr L_\bullet\otimes_{\mathscr O_Y}\mathscr M)$ pour
$i=i_1-1$. On voit donc qu'on peut supposer dans (7.7.4) que
$\mathscr L_\bullet$ est aussi un complexe \emph{fini}, à condition d'exiger
seulement que les $\mathscr L_i$ soient des $\mathscr O_Y$-Modules
\emph{localement libres} (associés à des $A$-modules projectifs de type fini).

Ce raisonnement s'applique en particulier au cas où $\mathscr P_\bullet$ est
réduit à \emph{un seul} terme $\mathscr F\ne0$, de degré 0 (auquel cas
$\mathscr T_n(\mathscr M)=\mathrm R^{-n}f_*
(\mathscr F\otimes_{\mathscr O_Y}\mathscr M)$) ; on peut alors supposer que les
$\mathscr L_i$ sont nuls pour $i>0$ ; on utilisera de préférence dans ce cas les
notations cohomologiques, écrivant donc $\mathscr T^{-p}$ au lieu de
$\mathscr T_p$.

\medskip
\textnormal{(ii)} Lorsque dans l'énoncé de (7.7.5), on ne suppose plus que les
$\mathscr P_i$ sont $Y$-plats, les conclusions restent valables à condition de
poser cette fois
\[
  \mathscr T_p(\mathscr M)
  =
  \boldsymbol{\mathscr{T}\!or}^Y_p
    (f,1_Y;\mathscr P_\bullet,\mathscr M).
  \tag{7.7.12.1}\label{III.7.7.12.1.fr}
\]
En effet, $\boldsymbol{\mathscr{T}\!or}^Y_n(f,1_Y;\mathscr P_\bullet,\mathscr O_Y)$ est
alors un $\mathscr O_Y$-Module \emph{cohérent} en vertu de (6.7.9). La
démonstration de (6.10.5) s'applique sans changement, compte tenu de (6.10.1)
et montre encore que lorsque $Y=\operatorname{Spec}(A)$ est affine, l'on a
$\mathscr T_p(\mathscr M)=\mathscr H_p(\mathscr L_\bullet\otimes_{\mathscr O_Y}\mathscr M)$,
avec $\mathscr L_\bullet=\widetilde{L}_\bullet$,
où $L_\bullet$ est un complexe de $A$-modules libres de type fini ; cela prouve
notre assertion.
\end{remarks}

\subsection{Application aux morphismes propres : II. Critères de platitude
cohomologique.}
\label{subsection:III.7.8.fr}

\begin{definition}[7.8.1]
Soient $X$, $Y$ deux préschémas, $f:X\to Y$ un morphisme quasi-compact et
séparé, $\mathscr P_\bullet$ un complexe de $\mathscr O_X$-Modules
quasi-cohérents et $Y$-plats, $\mathscr T_\bullet$ le foncteur homologique de
$\mathscr O_Y$-Modules quasi-cohérents défini par \textup{(7.7.1.1)}, $y$ un point de $Y$.
On dit que $\mathscr P_\bullet$ est
